\section{Concepts: the design tree, placement, and what persists}\label{sec:concepts}

\subsection{A rocket is a tree of parts}\label{sec:parts}

A QtRocket design is a tree structure: one root part (usually the nose cone) with
children attached to it, children attached to those, and so on.
\repl{newdesign} creates the root; \repl{addpart} attaches a child to any
existing part, addressed by the \code{name=} you gave it, by its numeric id,
or by the word \code{root}. Composite mass, center of gravity, and inertia are
computed from the tree --- there is no separate ``mass budget'' to maintain.

\begin{tip}[Address parts by name in scripts]
Ids are assigned in session order and \emph{reset when a design is reloaded},
so a saved script that says \code{removepart 5} may bite a different part
tomorrow. Names are stable: give every part a \code{name=} and use it.
\end{tip}

Four part types are constructible today. All lengths are meters, densities
kg/m$^3$ --- never millimeters.

\begin{table}[htbp]\centering\small
\begin{tabular}{@{}lll@{}}
\toprule
\textbf{Type} & \textbf{Required keys} & \textbf{Optional keys}\\
\midrule
\code{NoseCone}     & \code{baseRadius length density} & \code{solid} (default \code{true}), \code{wallThickness}, \code{name}\\
\code{BodyTube}     & \code{outerRadius length density} & \code{innerRadius} (default 0 $=$ solid rod), \code{name}\\
\code{FinSet}       & \code{finCount rootChord tipChord span} & \code{sweep} (default 0), \code{name}\\
                    & \code{thickness bodyRadius density} & \\
\code{HollowSphere} & \code{outerRadius density} & \code{innerRadius} (default 0 $=$ solid), \code{name}\\
\bottomrule
\end{tabular}
\caption{Constructible part types and their \code{key=value} parameters. A
motor is \emph{not} a factory part --- it attaches through \repl{setmotor}.}
\label{tab:parts}
\end{table}

Parameter values are parsed strictly: \code{length=0.25abc} is rejected, not
truncated. Unknown keys are skipped so newer files keep working on older
builds --- but a token without \code{=} is an error.

\subsection{Placement: seat, stations, gap}\label{sec:placement}

Every \repl{addpart} link records \emph{physical intent}, not
coordinates: a \textbf{seat} kind plus optional station and gap overrides.
Stations are fractions of a part's live length --- \code{0} is the aft plane,
\code{1} the fore plane --- so editing a part's length later moves every seam
with it.

\begin{table}[htbp]\centering\small
\begin{tabularx}{\linewidth}{@{}lXl@{}}
\toprule
\textbf{Seat} & \textbf{Meaning} & \textbf{\code{gap} means}\\
\midrule
\seat{Abut}       & Child's fore plane against the parent's aft plane --- how a
                    body tube hangs under a nose cone. The default seat. &
                    signed forward standoff ($0=$ touching)\\
\seat{NestInBore} & Child slides into the parent's bore from the fore rim ---
                    couplers, payloads. & insertion depth ($\ge 0$;
                    $=$ child length $\Rightarrow$ flush)\\
\seat{OnSurface}  & Child sits on the parent's outer wall, e.g. fins. Align the
                    axial position with \code{parentStation}/\code{childStation}. &
                    signed forward standoff along $+z$, like \seat{Abut}\\
\bottomrule
\end{tabularx}
\caption{The three seat kinds. The seat sets a canonical station pair first;
explicit \code{parentStation}/\code{childStation}/\code{gap} keys then override
individually.}
\label{tab:seats}
\end{table}

\repl{checkdesign} is the arbiter of whether the assembled geometry makes
physical sense: it reports \emph{overlaps} (a part intruding into another's
solid material) with the offending part ids, the radial penetration, and the
$z$ location, and \emph{air gaps} (a stack that is not contiguous) with the
uncovered $z$ interval and its size.
Tutorial~3 breaks a design twice on purpose to show both messages ---
and shows the visualizer painting the offender red.

\begin{note}[Further Reading]
The placement model --- stations, seats, the resolver, and its diagnostics ---
has its own design whitepaper in
\srcf{docs/PartPlacementDesignLatex/}. This guide just teaches the usage.
See the whitepaper for more detail.
\end{note}

\subsection{Motors and the motor database}\label{sec:motors-concept}

Motors live in a session database, filled from three sources that merge: local
RockSim files (\repl{loadmotors} \code{file.rse}), local RASP files
(\repl{loadmotors} \code{file.eng}), and QtRocket's own \code{.qmd} format
(\repl{loaddb}/\repl{savedb}). An online search
(\repl{tcsearch}, needs network) adds results from
\href{https://www.thrustcurve.org}{thrustcurve.org} into the same database.
\repl{setmotor} selects a motor by its exact common name (\code{B6},
\code{G80T}); \repl{listmotors} shows what is loaded.

\subsection{What persists, and where}\label{sec:persistence}

\begin{itemize}
\item \textbf{The design} --- part tree, placement links, drag coefficient,
      reference area, and the selected motor \emph{by name} --- is written by
      \repl{savedesign} to a versioned \code{.qrd} XML file and restored by
      \repl{loaddesign}. If the named motor is not in the database at load
      time, the design still loads without one, and the \ok{} line says so:
      \code{OK loaddesign: file.qrd (no motor)}.
\item \textbf{The motor database} persists only via \repl{savedb} to
      \code{.qmd}.
\item \textbf{Session configuration} --- atmosphere, gravity, integrator,
      timestep, launch velocity and angle --- is deliberately \emph{not} part
      of the design file. Scripts restate it; that keeps a \code{.qrd}
      a description of a rocket, not of one particular flight.
\end{itemize}
